\chapter{Anion Gap}

\section*{What Is This Test?}
The anion gap (AG) is a calculated value derived from the electrolyte panel that represents the difference between the measured serum cations (positively charged ions) and measured serum anions (negatively charged ions). The standard formula is: Anion Gap = [Na\textsuperscript{+}] -- ([Cl\textsuperscript{-}] + [HCO\textsubscript{3}\textsuperscript{-}]). The normal anion gap is 8--12 mEq/L (mmol/L), with some laboratories using a reference of 7--16 mEq/L depending on instrumentation. In a healthy person, the measured cation (sodium) exceeds the sum of the measured anions (chloride and bicarbonate) because normal plasma contains unmeasured anions that account for the gap: albumin (the dominant unmeasured anion, contributing ~2.5 mEq/L per 1 g/dL of albumin), phosphate (contributes ~2 mEq/L at normal phosphate), sulfate (~1 mEq/L), lactate (<2 mEq/L), and other organic acids. There are also unmeasured cations (potassium, calcium, magnesium) that are typically excluded from the standard formula. The anion gap is the single most useful laboratory parameter for classifying metabolic acidosis into two broad categories: (1) high-anion-gap metabolic acidosis, caused by accumulation of unmeasured acid anions (lactate, ketoacids, toxic alcohols, uremic acids), and (2) normal-anion-gap (hyperchloremic) metabolic acidosis, caused by loss of bicarbonate from the GI tract or kidneys with compensatory chloride retention. Because albumin is the major component of the normal anion gap, the gap must be corrected for hypoalbuminemia, which is common in hospitalized and critically ill patients. The albumin-corrected anion gap = measured AG + 2.5 $\times$ (4.0 -- measured albumin in g/dL).

\section*{Alternative Names}
AG; Anion Gap, Serum; Serum Anion Gap; Calculated Anion Gap; Delta Gap; Albumin-Corrected Anion Gap; Anion Gap (Na -- Cl -- HCO3)

\section*{Principle}
The anion gap is not a directly measured test but a mathematical calculation performed by the laboratory information system using the measured concentrations of sodium, chloride, and total CO\textsubscript{2} (bicarbonate). Sodium is measured by ion-selective electrode (ISE) potentiometry with a glass membrane electrode; chloride is measured by ISE with a quaternary ammonium chloride-sensitive membrane; bicarbonate is measured enzymatically via the PEPC method. The anion gap = [Na\textsuperscript{+}] -- ([Cl\textsuperscript{-}] + [HCO\textsubscript{3}\textsuperscript{-}]). Some formulas include potassium: AG = ([Na\textsuperscript{+}] + [K\textsuperscript{+}]) -- ([Cl\textsuperscript{-}] + [HCO\textsubscript{3}\textsuperscript{-}]), which yields a normal range of 12--16 mEq/L (the higher range reflects the inclusion of the ~4 mEq/L contribution from potassium). Most US laboratories use the formula without potassium (Na -- Cl -- HCO\textsubscript{3}), with a normal range of 8--12 mEq/L. The calculation is performed automatically by the chemistry analyzer or laboratory information system and appears on the patient's report alongside the electrolyte values. Accuracy of the anion gap depends entirely on the accuracy of the three measured components---errors in sodium, chloride, or bicarbonate measurement will propagate to the anion gap. Notably, a spuriously elevated chloride (from bromide interference on the chloride ISE) will produce a falsely low or negative anion gap, while a spuriously low bicarbonate (from uncapped tube CO\textsubscript{2} off-gassing) will produce a falsely elevated anion gap.

\section*{History}
The concept of the anion gap was first described by James Gamble in the 1940s in his diagrams of electrolyte composition (Gamblegrams). However, the clinical utility of the anion gap in diagnosing metabolic acidosis was recognized and popularized in the 1970s. The seminal paper by Emmet and Narins (1977) established the modern clinical approach to metabolic acidosis using the anion gap, with the mnemonic MUDPILES for high-AG acidosis later adapted by subsequent authors \cite{emmett1977clinical}. Correction for albumin was proposed in the 1990s when it became clear that hypoalbuminemia (ubiquitous in hospitalized and ICU patients) masks true elevations of the anion gap. Figge and colleagues (1998) published the widely adopted albumin correction factor: for every 1 g/dL decrease in albumin below 4.0, add 2.5 mEq/L to the calculated anion gap \cite{figge1998anion}. The delta-delta (delta gap) approach to detecting mixed acid-base disorders was developed by Wrenn and colleagues in the 1990s \cite{wrenn1990delta}. Modern laboratory information systems now calculate and display the anion gap automatically with every BMP and CMP.

\section*{Why This Test Is Ordered}
\begin{itemize}
  \item Classification of metabolic acidosis: the anion gap is the essential first step in determining whether a metabolic acidosis is a high-AG acidosis (accumulation of unmeasured acid anions) or a normal-AG (hyperchloremic) acidosis (bicarbonate loss). This distinction drastically narrows the differential diagnosis
  \item Workup of unexplained low bicarbonate (<22 mEq/L) on a routine electrolyte panel: any low bicarbonate should prompt anion gap calculation and classification of the acidosis
  \item Detection of toxic alcohol ingestion (methanol, ethylene glycol): high-AG metabolic acidosis with elevated osmolal gap in a patient with altered mental status or suspected ingestion
  \item Monitoring of diabetic ketoacidosis (DKA): the anion gap is used to track resolution of ketoacidosis during insulin and fluid therapy---closure of the anion gap (return to 8--12) indicates that ketosis has resolved, even if the patient still has hyperchloremic acidosis from normal saline resuscitation
  \item Detection of lactic acidosis in sepsis, shock, cardiac arrest, mesenteric ischemia, and metformin toxicity: the anion gap rises proportionally to lactate accumulation (approximately 1:1 in early lactic acidosis)
  \item Identification of mixed acid-base disorders using delta-delta (delta gap) analysis when both high-AG metabolic acidosis and another acid-base disorder are suspected
  \item Investigation of unexplained low or negative anion gap: suggests hypoalbuminemia, laboratory error, bromide toxicity, lithium toxicity, or IgG paraproteinemia (multiple myeloma)
  \item Monitoring of salicylate poisoning: aspirin causes a mixed high-AG metabolic acidosis and primary respiratory alkalosis (from direct stimulation of medullary respiratory center)
  \item Uremic acidosis monitoring in advanced chronic kidney disease (GFR <15--20 mL/min)
\end{itemize}

\section*{Primary vs Secondary Test}
The anion gap is a secondary calculated parameter. It is derived automatically from the electrolyte panel results and is not a separately orderable test. However, it is one of the most powerful and frequently used calculated values in clinical medicine for the differential diagnosis of metabolic acidosis. Any bicarbonate <22 mEq/L should immediately prompt anion gap calculation and classification. The anion gap is automatically reported with every BMP and CMP in most laboratory information systems.

\section*{How This Test Is Performed}
No additional blood draw is required beyond the standard venipuncture for the basic metabolic panel. The anion gap is automatically calculated by the laboratory information system using the measured sodium, chloride, and bicarbonate values from the same blood sample collected in a serum separator tube (gold-top) or lithium heparin tube (green-top). The formula applied is: Na -- (Cl + HCO\textsubscript{3}). The result appears on the patient report alongside the other electrolyte values. No separate test ordering is required; the anion gap is a derived value that reports with every BMP or CMP unless laboratories suppress it. Because the anion gap depends entirely on the accuracy of the three measured analytes, proper sample handling is critical: capped tubes to prevent CO\textsubscript{2} loss (which lowers bicarbonate and falsely elevates the AG), avoidance of hemolysis, prompt centrifugation, and avoidance of prolonged tourniquet time.

\section*{What Preparation Is Needed}
No special patient preparation beyond that for the basic metabolic panel. The same pre-analytical precautions apply: (1) The blood tube must be filled completely and kept capped until centrifugation---CO\textsubscript{2} escape to room air falsely lowers bicarbonate, causing a falsely elevated anion gap. (2) Prolonged tourniquet time should be avoided (hemoconcentration may slightly alter electrolytes). (3) Blood should not be drawn from an arm receiving IV fluids. (4) No fasting is required for the anion gap itself, though fasting may be necessary if the sample is being used for other tests (glucose, lipids) as part of a CMP. (5) Inform the healthcare provider of medications that affect acid-base status: metformin (lactic acidosis risk), aspirin, diuretics, sodium bicarbonate, acetazolamide, topiramate, and corticosteroids. (6) If the patient has received large-volume normal saline resuscitation, the chloride load may cause a superimposed hyperchloremic metabolic acidosis, lowering the anion gap even while a high-AG acidosis is resolving.

\section*{Contraindications and Precautions}
No absolute contraindications. The reliability of the calculated anion gap depends on the accuracy of the component measurements. Important considerations: (1) Hypoalbuminemia: this is the most important correction because critically ill patients almost universally have low albumin. For every 1 g/dL decrease in albumin below 4.0 g/dL, the expected anion gap decreases by 2.5 mEq/L. A patient with albumin of 2.0 g/dL has an expected pre-pathology anion gap of only 7 mEq/L ([12] -- [2.5 $\times$ 2] = 7). If the measured AG is 12 in this patient (seemingly "normal"), the albumin-corrected AG is actually 17 mEq/L---a significant high-AG metabolic acidosis. Failure to apply albumin correction is the most common error in anion gap interpretation. (2) Bromide toxicity: bromide interferes with the chloride ISE (selectivity coefficient ~2.5:1), causing pseudohyperchloremia. This produces an apparently elevated chloride, causing a falsely low or negative anion gap with no true acid-base disturbance. A negative anion gap should always prompt consideration of bromide toxicity. (3) Lithium: lithium is a small cation that is not included in the anion gap calculation, so severe lithium toxicity can create a slightly low anion gap. (4) Multiple myeloma (IgG paraprotein): the IgG M-protein has a net positive charge at physiological pH (isoelectric point 7.3--8.0), behaving as a cation and lowering the true anion gap. This can result in a low or even negative anion gap. In contrast, IgA paraprotein has a net negative charge and may slightly elevate the anion gap. (5) Hyperlipidemia: severe hypertriglyceridemia can cause pseudohyponatremia with indirect ISE, lowering the measured sodium and thus the calculated anion gap. (6) Delayed bicarbonate measurement or uncapped tubes: CO\textsubscript{2} loss lowers bicarbonate, falsely elevating the calculated anion gap. If the sample was uncapped for a prolonged period, the anion gap should be interpreted with caution. (7) Laboratory error in any of the three component measurements (sodium, chloride, bicarbonate) will propagate directly into the anion gap.

\section*{Risks}
The anion gap is a calculated value with no direct risk to the patient. Risks are limited to those of the venipuncture required to obtain the BMP: mild pain, bruising, hematoma, vasovagal syncope (1--3\% of draws), and rare infection or nerve injury. The blood volume drawn (3--5 mL) is negligible.

\section*{What Happens After the Test}
The anion gap result is reported automatically as part of the electrolyte panel, typically within 30--60 minutes. The clinician's interpretation follows a stepwise approach: (1) Is the bicarbonate low (<22 mEq/L)? If yes, a metabolic acidosis is present, and the anion gap must be calculated and albumin-corrected. (2) Is the albumin-corrected AG elevated (>12 mEq/L)? If yes: high-AG metabolic acidosis. If no: normal-AG (hyperchloremic) metabolic acidosis. (3) For high-AG acidosis, determine etiology: check serum lactate, beta-hydroxybutyrate/ketones, toxic alcohol screen, salicylate level, renal function. Calculate osmolal gap if toxic ingestion suspected. (4) Perform delta-delta analysis: (measured AG -- 12) + measured HCO\textsubscript{3}. If 24--28: pure high-AG metabolic acidosis. If >28: concurrent metabolic alkalosis is also present (e.g., vomiting with DKA). If <22: concurrent normal-AG metabolic acidosis is also present (e.g., diarrhea with DKA). (5) For normal-AG acidosis, calculate urine anion gap [urine Na + urine K -- urine Cl] to differentiate GI bicarbonate loss (negative urine AG, reflecting NH\textsubscript{4}Cl excretion) from renal tubular acidosis (positive urine AG). (6) Serial anion gap monitoring is critical in DKA treatment---closure of the gap to normal indicates resolution of ketoacidosis; persistent elevation indicates ongoing ketone production requiring continued insulin. (7) If anion gap is negative or very low, investigate for bromide, laboratory error, hypoalbuminemia, lithium, or IgG myeloma.

\section*{Understanding Results}

\subsection*{Below Normal}
\begin{itemize}
  \item \textbf{Low anion gap (<3 mEq/L after albumin correction):} Uncommon but clinically significant. In a patient with normal albumin, the anion gap should never normally be <3 mEq/L.
  \item \textbf{Hypoalbuminemia (most common cause):} Albumin is the dominant unmeasured anion. For every 1 g/dL decrease in albumin below 4.0 g/dL, the expected anion gap decreases by 2.5 mEq/L. Thus, an albumin of 2.0 g/dL produces an expected baseline AG of only 7 mEq/L. This is not a true "low" AG but requires albumin correction for accurate interpretation. Critically ill patients, those with cirrhosis, nephrotic syndrome, malnutrition, and burns commonly have hypoalbuminemia.
  \item \textbf{Laboratory error:} Underestimated sodium (pseudohyponatremia from hypertriglyceridemia or hyperproteinemia with indirect ISE) overestimates the unmeasured cation portion and produces a spuriously low AG. Overestimated chloride (bromide interference) is the most common lab artifact causing a low or negative AG. Underestimated bicarbonate (CO\textsubscript{2} gas loss from uncapped tube) is unlikely since this would elevate the AG.
  \item \textbf{Bromide toxicity:} Bromide (Br\textsuperscript{-}) is sensed by the chloride ISE with approximately 2.5$\times$ greater affinity than chloride itself. The analyzer reports this as chloride, producing pseudohyperchloremia with apparent elevated chloride and consequently a low or negative anion gap. Historical bromism from over-the-counter sedatives (e.g., Bromo-Seltzer, now rare), but also from pyridostigmine bromide (used in myasthenia gravis, though doses are usually too low) and some herbal or imported medications.
  \item \textbf{Lithium toxicity:} Lithium (Li\textsuperscript{+}) is a small monovalent cation not included in the anion gap formula. At toxic levels (>1.5 mmol/L), lithium contributes to the unmeasured cation pool, producing a slightly low anion gap. This is rarely the sole diagnostic clue but may be seen when lithium levels exceed 3--4 mmol/L.
  \item \textbf{Multiple myeloma (IgG paraprotein):} IgG immunoglobulin has a net positive charge at physiological pH (isoelectric point 7.3--8.0), acting as an unmeasured cation. A large IgG M-protein burden can produce a low or even negative anion gap. IgA paraproteins are negatively charged and may produce a slightly elevated AG. This is an important diagnostic clue: an unexplained low or negative AG in an adult >50 years should prompt serum protein electrophoresis (SPEP) to evaluate for monoclonal gammopathy.
  \item \textbf{Hypercalcemia, hypermagnesemia, hyperkalemia:} These cations are excluded from the standard AG formula. While normal concentrations contribute negligibly, extreme elevations (>14 mg/dL calcium, >5 mEq/L magnesium, >7 mEq/L potassium) may theoretically lower the AG slightly, though this is rarely clinically significant.
  \item \textbf{Hyponatremia:} Severe hyponatremia (<125 mEq/L) reduces the numerator of the AG calculation, potentially producing a lower absolute AG value. This is not a true acid-base abnormality.
\end{itemize}

\subsection*{Above Normal}
\begin{itemize}
  \item \textbf{Elevated anion gap (>12 mEq/L with albumin correction):} Indicates high-AG metabolic acidosis due to accumulation of unmeasured acid anions in the blood. Mnemonic for causes: MUDPILES.
  \item \textbf{Methanol (M):} Metabolized by alcohol dehydrogenase and aldehyde dehydrogenase to formic acid, a potent mitochondrial toxin. Causes high-AG metabolic acidosis with elevated osmolal gap (>10 mOsm/kg), visual disturbances (blurred vision, photophobia, "snowfield vision"), retinal toxicity, and potential blindness. Often from windshield washer fluid, moonshine, or industrial solvents. Treatment: fomepizole (alcohol dehydrogenase inhibitor) or ethanol, plus hemodialysis if levels elevated or end-organ damage present.
  \item \textbf{Uremia (U):} Advanced chronic kidney disease (GFR <15--20 mL/min) impairs excretion of organic acid anions (sulfate, phosphate, urate, hippurate). Accumulation of these unmeasured anions increases the AG. The AG in uremic acidosis rarely exceeds 20--24 mEq/L; a higher AG suggests superimposed lactic acidosis or another cause. Treatment: renal replacement therapy (hemodialysis).
  \item \textbf{Diabetic ketoacidosis (D):} Accumulation of acetoacetate and beta-hydroxybutyrate (ketoanions). AG >20 mEq/L is common; >24 mEq/L correlates with severe ketoacidosis. DKA is the most common cause of high-AG metabolic acidosis in patients under 40 years. Predominant ketone is beta-hydroxybutyrate (not detected by urine ketone dipsticks). Closure of the AG (return to 8--12 mEq/L) signals resolution of ketoacidosis. Treatment: IV insulin, aggressive fluid resuscitation, and potassium repletion.
  \item \textbf{Propylene glycol (P):} Used as a solvent for IV medications (lorazepam, diazepam, phenytoin, esmolol, nitroglycerin). Metabolized to lactic acid, causing high-AG metabolic acidosis with elevated lactate and osmolal gap. Typically occurs with prolonged high-dose IV infusions in ICU settings. Treatment: discontinue the offending agent; supportive care.
  \item \textbf{Isoniazid/Iron (I):} Isoniazid (INH) overdose causes metabolic acidosis through multiple mechanisms, including generation of lactic acid from seizure activity (INH causes pyridoxine deficiency $\rightarrow$ inhibits GABA synthesis $\rightarrow$ refractory seizures $\rightarrow$ lactic acidosis). Also may form isonicotinic acid directly. Treatment: pyridoxine (vitamin B6) for seizures, benzodiazepines, supportive care. Iron overdose causes direct mitochondrial toxicity with lactic acidosis, GI hemorrhage, and hepatotoxicity.
  \item \textbf{Lactic acidosis (L):} Most common cause of elevated AG. Type A (tissue hypoxia): shock (septic, cardiogenic, hypovolemic), cardiac arrest, severe anemia, mesenteric ischemia, carbon monoxide poisoning, cyanide poisoning, severe hypoxemia. Type B (impaired cellular metabolism without global hypoxia): sepsis (mitochondrial dysfunction), malignancy (Warburg effect), metformin toxicity, liver failure, mitochondrial disorders, thiamine deficiency, D-lactate from short bowel syndrome with bacterial overgrowth. Lactate >4 mmol/L is associated with significantly increased mortality in sepsis (28--38\%). Treatment: address underlying cause (resuscitation, antibiotics, source control for sepsis; discontinue metformin; thiamine repletion).
  \item \textbf{Ethylene glycol (E):} Found in antifreeze, brake fluid, de-icing solutions. Metabolized to glycolic acid and oxalic acid. Causes high-AG metabolic acidosis with elevated osmolal gap, calcium oxalate crystalluria (envelope-shaped crystals), acute kidney injury, and eventual cardiopulmonary failure. The elevated AG combined with elevated osmolal gap is classic. Treatment: fomepizole (preferred) or ethanol, plus urgent hemodialysis.
  \item \textbf{Salicylates (S):} Aspirin poisoning causes a mixed acid-base disorder: (a) direct stimulation of the medullary respiratory center produces primary respiratory alkalosis; (b) uncoupling of oxidative phosphorylation produces metabolic acidosis with elevated lactate and ketone bodies. In adults, the classic picture is mixed high-AG metabolic acidosis plus respiratory alkalosis. In children, metabolic acidosis predominates. Treatment: alkalinization of urine with sodium bicarbonate (enhances salicylate excretion), hemodialysis for severe toxicity (>100 mg/dL acute, >60 mg/dL chronic with end-organ damage).
  \item \textbf{Additional causes not in MUDPILES:} (a) Pyroglutamic acidosis (5-oxoprolinemia): from chronic acetaminophen use in setting of glutathione depletion (malnutrition, sepsis, chronic illness); causes elevated AG with elevated 5-oxoproline levels in urine. (b) Alcoholic ketoacidosis: occurs in chronic alcoholics with binge drinking followed by vomiting, decreased food intake; presents with high-AG metabolic acidosis, elevated beta-hydroxybutyrate (predominant), low or normal glucose. (c) Rhabdomyolysis: release of organic acids and phosphate from damaged muscle, with acute kidney injury compounding the acidosis.
  \item \textbf{Delta-delta (delta gap) analysis for mixed disorders:} Formula: (measured AG -- 12) + measured HCO\textsubscript{3}. Sum should approximate 24--28 mEq/L in a pure high-AG metabolic acidosis. If sum >28: concurrent metabolic alkalosis present (e.g., vomiting patient with DKA, diuretic use in patient with lactic acidosis). If sum <22: concurrent normal-AG (hyperchloremic) metabolic acidosis present (e.g., diarrhea plus DKA, renal tubular acidosis plus lactic acidosis) \cite{wrenn1990delta}.
  \item \textbf{Anion gap >20 mEq/L:} Strongly associated with DKA, toxic alcohols, or severe lactic acidosis. Requires aggressive investigation and treatment. Anion gap >25 mEq/L is associated with significant morbidity and mortality.
\end{itemize}

\section*{Normal Results}
\begin{itemize}
  \item \textbf{Standard formula (Na -- Cl -- HCO\textsubscript{3}):} 8--12 mEq/L
  \item \textbf{With potassium [Na + K -- Cl -- HCO\textsubscript{3}]:} 12--16 mEq/L
  \item \textbf{Albumin-corrected anion gap:} 8--12 mEq/L. Correction: AG\textsubscript{corrected} = AG\textsubscript{measured} + 2.5 $\times$ (4.0 -- albumin in g/dL)
  \item \textbf{Expected anion gap at specific albumin levels:} Albumin 4.0 g/dL $\rightarrow$ AG 12; Albumin 3.0 g/dL $\rightarrow$ AG 9.5; Albumin 2.0 g/dL $\rightarrow$ AG 7; Albumin 1.0 g/dL $\rightarrow$ AG 4.5
  \item \textbf{Critical finding:} Negative anion gap (should always be investigated---consider bromide, IgG myeloma, laboratory error)
\end{itemize}

\section*{Follow-up Tests}
\begin{itemize}
  \item \textbf{Serum albumin:} Essential for calculating albumin-corrected anion gap. Should be ordered from the same sample as the electrolyte panel for accurate correction.
  \item \textbf{Serum lactate:} First test to order for elevated AG. Most common cause of high-AG metabolic acidosis. Normal <2 mmol/L. Point-of-care lactate available in most emergency departments.
  \item \textbf{Serum beta-hydroxybutyrate:} For suspected DKA or alcoholic ketoacidosis. Preferred over urine ketones (which do not detect beta-hydroxybutyrate, the dominant ketone in DKA). >3 mmol/L indicates significant ketosis.
  \item \textbf{Serum osmolality (measured) and osmolal gap:} Calculated osmolality = 2(Na) + Glucose/18 + BUN/2.8 + Ethanol/4.6. Osmolal gap = measured -- calculated. Normal <10 mOsm/kg. Gap >10 mOsm/kg with elevated AG suggests methanol or ethylene glycol intoxication.
  \item \textbf{Toxic alcohol screen (methanol, ethylene glycol, isopropanol):} Order urgently if high-AG metabolic acidosis with elevated osmolal gap. Isopropanol does NOT cause high-AG acidosis (it is metabolized to acetone, a ketone, without acidosis) but contributes to osmolal gap.
  \item \textbf{Serum salicylate level:} If salicylate poisoning suspected. Therapeutic range: 10--30 mg/dL. Toxicity: >30 mg/dL. Severe toxicity: >100 mg/dL (acute), >60 mg/dL (chronic).
  \item \textbf{Urine calcium oxalate crystals:} Microscopic examination of urine for envelope-shaped (dihydrate) or needle-shaped (monohydrate) calcium oxalate crystals, highly suggestive of ethylene glycol poisoning. Also check for urine fluorescence under Wood's lamp (ethylene glycol in some antifreeze formulations contains fluorescein).
  \item \textbf{Serum iron level:} If iron overdose is suspected (GI hemorrhage, hepatotoxicity, high-AG metabolic acidosis).
  \item \textbf{Serum protein electrophoresis (SPEP) and immunofixation:} If low or negative AG suggests IgG multiple myeloma (unmeasured cationic paraprotein).
  \item \textbf{Serum bromide level:} If pseudohyperchloremia with negative AG suggests bromide toxicity.
  \item \textbf{Urine anion gap:} Calculated as urine Na + urine K -- urine Cl. Used in normal-AG (hyperchloremic) metabolic acidosis to differentiate GI bicarbonate loss (negative urine AG, indicating appropriate NH\textsubscript{4}Cl excretion) from renal tubular acidosis (positive urine AG, indicating impaired renal acidification).
\end{itemize}

\section*{References}
\bibliographystyle{unsrtnat}
\bibliography{references}

\clearpage
\section*{Regulatory Status and Authorization Date}
Regulatory status applies to the specific assay, instrument, manufacturer, and intended use rather than to the test name alone. No single approval date can be assigned to this test category. For a commercial assay, verify the current product in the relevant regulator's database: the U.S. Food and Drug Administration (FDA) clearance, approval, or authorization; India's Central Drugs Standard Control Organisation (CDSCO) licence or permission; the European Union In Vitro Diagnostic Regulation (EU IVDR) CE certificate issued through the applicable conformity-assessment route; or the United Kingdom Medicines and Healthcare products Regulatory Agency (MHRA) registration and UKCA/CE status. Laboratory-developed tests may not have an individual FDA, CDSCO, EU IVDR, or MHRA approval date.
\clearpage
